\chapter{总结}\label{ux7b2cux516dux7ae0-ux603bux7ed3}

文本到语音合成、语音克隆和大语言模型驱动的生成式语音技术快速发展，使真实语音被未授权采集、编码和复刻的门槛不断降低。攻击者可利用公开视频、直播录屏、有声读物、企业客服语音等渠道获取目标说话人的声音样本，通过零样本语音克隆、少样本微调式语音克隆或基于语音分词器的语音合成系统，生成高度逼真的伪造语音，进而引发身份冒用、电信诈骗、虚假内容传播和个人声音权益滥用等安全风险。声音作为一种具有身份属性、传播频繁且难以更换的数字资产，已成为生成式人工智能安全治理中的重要保护对象。现有事后检测手段难以阻止攻击者在语音公开发布后提前采集和建模，亟需在数据源头建立主动防护能力。

本作品设计并实现了一套面向声音资产发布前防护的主动式语音克隆防护系统V-Guard。系统以语义表示和声音身份表示的联合扰动为核心，在发布前对原始音频施加受心理声学约束的保护扰动，使公开语音在保持正常听感的同时降低对未授权克隆系统的可利用价值，并将上述机制封装为完整的网页端闭环平台。具体而言，本作品的贡献主要体现在以下几个方面。

第一，本作品提出了一种在发布源头实施保护的主动式防克隆方案。该方案将防护位置从事后检测前移至语音公开发布之前，通过为原始音频施加保护扰动，使公开语音在保持正常听感的同时降低对未授权克隆系统的可利用价值。方案不依赖于对攻击者所用克隆模型的先验知识，能够覆盖零样本、少样本和微调式等多种克隆方式。为支撑上述方案的落地使用，本作品开发了完整的网页端系统，提供音频或视频上传、自动预处理、保护扰动生成、ASR转写对比、说话人相似度评估、克隆防护验证、音频质量分析、历史任务管理以及报告导出等功能，为个人用户、内容创作者和企业声音资产管理方提供了一种可直接使用的主动防护工具。

第二，在技术机制上，V-Guard在技术实现上采用了语义-身份联合防护的设计路线。语义链路将防护目标从ASR输出文本层面下沉至语音编码器的中间表示层，利用离散语音token量化过程的敏感性引发token序列的结构性变化，从而破坏下游模型对语音内容的理解；身份链路通过集成多种编码器覆盖不同模型对声音身份的提取方式，使保护扰动在多个身份表示空间中同时偏离原始说话人，避免对单一模型过拟合。心理声学掩蔽模型的引入则确保扰动在有效干扰机器模型的同时不会明显影响人耳听感。核心算法采用多目标优化结构，通过代理模型组近似建模语音克隆系统的处理过程，冻结模型参数并仅优化输入侧保护扰动，兼顾连续语音表示与离散语音
token 两种建模路径。综合损失函数为：

\[\mathcal{L}\  = \ \lambda_{sem}\mathcal{L}_{sem}\  + \ \lambda_{id}\mathcal{L}_{id}\  + \ \lambda_{psy}\mathcal{L}_{psy} + \ \lambda_{2}\mathcal{L}_{2}\]

上述技术在V-Guard中，确保了扰动的有效性，鲁棒性和可迁移性，实际效果也验证了在表示层实施语义干扰以及多编码器集成覆盖身份提取这两种思路在语音防克隆场景中的有效性。

第三，本作品通过实验验证了系统的实际防护效果，同时也展示了该系统作为一项工程化工具的真实能力。基于LibriTTS测试样本的评估结果表明，受保护音频在多个ASR模型上的语义识别错误率显著提升，以保护音频为参考的克隆语音在说话人相似度上明显下降，证明了所采用的技术方法在真实场景中的有效性。高保真参数组、跨克隆后端迁移性和信道变换鲁棒性等实验进一步表明，系统具备可调节的防护-保真权衡能力和跨模型迁移性，也说明该系统在实际部署中具有一定的适应性和稳定性。上述结果不仅验证了V-Guard在语义和身份双链路上的防护有效性，也反映出该系统在个人声音隐私保护、企业声音资产管理和反诈安全教育等场景中具有进一步发展的潜力和研究前景。

总体来看，围绕声音资产源头主动防护这一目标,V-Guard为声音资产发布环节提供了一道前置保护层，在现有事后检测和内容审核手段之外，提供了一种在时间上更早介入的防护选择。\protect\phantomsection\label{_Toc617022676}{}

\FloatBarrier
